\documentclass[11pt]{article}

\usepackage[margin=1in]{geometry}
\usepackage{amsmath}
\usepackage{amssymb}
\usepackage{array}
\usepackage{hyperref}
\usepackage{enumitem}

\title{Toy Robot Arm Webcam Feedback Control Theory}
\author{toy\_robot\_arm project}
\date{July 3, 2026}

\begin{document}

\maketitle

\section{Purpose}

This note spells out the intended control loop for webcam-guided object
localization and robot arm positioning. The current code already contains the
core interfaces for arm poses, camera-derived 3D points, a proportional visual
servo step, and conservative servo-frame encoding. The missing pieces are the
runtime perception pipeline, calibrated coordinate transforms, and eventually a
live serial transport.

The key architectural rule is that perception can be simple fiducials, classical
computer vision, or machine learning, but the controller should see the same
kind of measurement:

\[
    z_t = (\hat{x}_t, \hat{y}_t, \hat{z}_t, c_t, \sigma_t)
\]

where $\hat{x}_t,\hat{y}_t,\hat{z}_t$ are the estimated target or tool position
in robot/workspace coordinates, $c_t$ is detection confidence, and $\sigma_t$
summarizes uncertainty. That keeps ML replaceable and testable instead of
letting a model directly command servo pulses.

\section{Control Loop Overview}

One iteration of the loop should be:

\begin{enumerate}[leftmargin=*]
    \item Capture synchronized webcam frame(s).
    \item Detect and localize the gripper/tool, target object, and optional drop
          zone markers.
    \item Convert image measurements into workspace coordinates.
    \item Estimate or smooth the current positions and reject low-confidence
          measurements.
    \item Compute the position error between the goal and observed tool pose.
    \item Convert that Cartesian error into bounded servo pulse deltas.
    \item Clamp the resulting pose to calibrated joint limits.
    \item Encode and send a slow servo move frame.
    \item Repeat after motion and camera latency settle.
\end{enumerate}

In block form:

\[
\begin{array}{c}
\text{webcam frames}
\\ \downarrow
\\
\text{detector/localizer}
\\ \downarrow
\\
\text{pixel or keypoint observations}
\\ \downarrow
\\
\text{camera calibration and workspace transform}
\\ \downarrow
\\
\text{filtered tool/object/drop-zone measurements}
\\ \downarrow
\\
\text{visual servo controller}
\\ \downarrow
\\
\text{bounded servo pulse commands}
\\ \downarrow
\\
\text{servo controller and physical arm}
\end{array}
\]

\section{Coordinate Frames}

The loop needs explicit frames:

\begin{description}
    \item[Image frame $I$] Pixel coordinates $(u,v)$ from a webcam.
    \item[Camera frame $C$] Metric coordinates relative to the camera optical
          center.
    \item[Workspace frame $W$] Table or robot-base coordinates in meters.
    \item[Servo space $S$] PWM pulse commands for channel IDs.
\end{description}

The current Rust type \texttt{Point3} represents a metric point:

\[
    p = (x_m, y_m, z_m)
\]

For a rectified stereo pair, the existing \texttt{StereoRig} implements:

\[
    z = \frac{fB}{d}, \qquad
    x = \frac{(u_L - c_x)z}{f}, \qquad
    y = \frac{(v_L - c_y)z}{f}
\]

where $f$ is focal length in pixels, $B$ is baseline in meters, and
$d = u_L - u_R$ is horizontal disparity.

For the first tabletop demo, an overhead monocular camera is simpler. If object
motion is mostly on a known table plane, a calibrated homography $H$ maps image
pixels into workspace coordinates:

\[
    \lambda
    \begin{bmatrix}
    x_W \\ y_W \\ 1
    \end{bmatrix}
    =
    H
    \begin{bmatrix}
    u \\ v \\ 1
    \end{bmatrix}
\]

Height $z_W$ can then be fixed to the table, inferred from a side camera, or
estimated later by stereo.

\section{Perception Measurements}

The perception layer should produce typed observations rather than direct arm
commands:

\[
    y_t = \{\text{label}, p_W, q, \Sigma, t\}
\]

where \textit{label} is tool/object/drop-zone, $p_W$ is a workspace position,
$q$ is confidence, $\Sigma$ is an optional covariance or uncertainty estimate,
and $t$ is frame time.

\subsection{Fiducial First Path}

The safest first path is fiducials:

\begin{itemize}[leftmargin=*]
    \item A marker on a rigid flag near the gripper provides tool position.
    \item A marker on the object provides target position.
    \item Workspace corner markers calibrate the table homography.
    \item A drop-zone marker gives a fixed placement goal.
\end{itemize}

Fiducials give repeatable labels and geometry while the arm is still being
calibrated. They also provide ground truth data for training or validating a
later ML detector.

\subsection{ML Detection and Localization Path}

ML becomes useful when objects are not fiducial-tagged, when lighting changes,
or when the system needs class-specific grasp points. The model should still
fit behind the same measurement interface. Suitable outputs are:

\begin{itemize}[leftmargin=*]
    \item Bounding boxes for object detection.
    \item Segmentation masks for object extent and centroid.
    \item Keypoints for gripper tip, object handle, or graspable features.
    \item Dense depth or monocular pose estimates, if calibrated well enough.
\end{itemize}

A practical ML path is:

\begin{enumerate}[leftmargin=*]
    \item Start with fiducial detections and logged images.
    \item Label object masks, centroids, and grasp keypoints from those logs.
    \item Train or fine tune a small detector/segmenter.
    \item Compare ML measurements against fiducial/calibration measurements.
    \item Gate control on confidence and uncertainty before allowing motion.
\end{enumerate}

The controller should ignore or slow down on weak measurements:

\[
    \text{accept}(y_t) =
    q_t \ge q_{\min}
    \quad \land \quad
    \|\Sigma_t\| \le \Sigma_{\max}
    \quad \land \quad
    t_{\text{now}} - t_t \le \Delta t_{\max}
\]

\section{State Estimation and Filtering}

Raw detections will jitter. A minimal filter can be an exponential moving
average:

\[
    \bar{p}_t = \alpha p_t + (1-\alpha)\bar{p}_{t-1}
\]

with lower $\alpha$ for noisy measurements and higher $\alpha$ for responsive
tracking. Later, a constant-velocity Kalman filter can account for camera
latency and measurement uncertainty:

\[
    x_t =
    \begin{bmatrix}
    p_t \\ v_t
    \end{bmatrix}
\]

For this arm, filtering should be conservative. It is better to hold position
when measurements disagree than to chase noise with servo motion.

\section{Visual Servo Controller}

The current controller is position-based visual servoing. It does not yet solve
full inverse kinematics. It computes a measured Cartesian error:

\[
    e_t = p_{\text{goal},t} - p_{\text{tool},t}
\]

The Rust code currently maps this error into servo pulse deltas using diagonal
proportional gains:

\[
\begin{aligned}
    \Delta s_1 &= k_x e_x \\
    \Delta s_2 &= k_{z,shoulder} e_z \\
    \Delta s_3 &= k_{z,elbow} e_z \\
    \Delta s_4 &= k_y e_y
\end{aligned}
\]

These are deliberately crude gains. They are useful before link lengths and
joint axes are measured because they let us jog the arm in the direction that
reduces visual error.

The implemented step then:

\begin{enumerate}[leftmargin=*]
    \item Computes Euclidean error norm $\|e_t\|$.
    \item Holds the current pose if $\|e_t\|$ is within tolerance.
    \item Rounds each pulse delta to an integer.
    \item Rate-limits each delta to \texttt{max\_delta\_pulse}.
    \item Adds deltas to the current pose.
    \item Clamps each known joint to calibrated limits.
\end{enumerate}

In compact form:

\[
    s_{t+1} =
    \operatorname{clamp}_{limits}
    \left(
        s_t +
        \operatorname{sat}_{\Delta s_{\max}}(K e_t)
    \right)
\]

where $s_t$ is the pulse vector, $K$ is the gain mapping, and
$\operatorname{sat}$ applies the per-step rate limit.

\section{Goal Selection}

The goal depends on the task phase:

\begin{description}
    \item[Approach] Put the gripper above the object with a safe height offset.
    \item[Descend] Move down while keeping x/y error small.
    \item[Close] Command the gripper from open to closed pulse.
    \item[Lift] Raise the object before lateral motion.
    \item[Transfer] Move above the drop zone.
    \item[Release] Lower, open the gripper, and retreat.
\end{description}

This should be implemented as a state machine. Each state supplies a
\texttt{VisualServoGoal}, tolerance, speed limit, and allowed gripper command.
The feedback controller remains the same.

\section{Safety and Gating}

The arm should only move when all gates pass:

\begin{itemize}[leftmargin=*]
    \item Servo IDs and joint pulse limits are calibrated.
    \item Current pose is known or initialized near neutral.
    \item Camera calibration is loaded and recent.
    \item Tool and target measurements pass confidence and freshness gates.
    \item Requested pulse deltas are under the rate limit.
    \item Resulting pose is inside joint limits.
    \item The task state permits the command.
\end{itemize}

If any gate fails, the controller should hold pose or return to a known safe
state. It should not extrapolate for long periods without camera feedback.

\section{Code Mapping}

\begin{tabular}{>{\ttfamily}p{0.33\linewidth}p{0.57\linewidth}}
\textnormal{Module} & Role \\
\hline
src/vision.rs & Camera geometry types and stereo triangulation helpers. \\
src/control.rs & Proportional visual servo step from measured error to pulse deltas. \\
src/arm.rs & Joint IDs, limits, neutral pose, validation, and clamping. \\
src/transport.rs & Servo controller frame encoding and transport abstraction. \\
src/main.rs & CLI simulation and frame inspection. \\
\end{tabular}

\section{Implementation Roadmap}

\begin{enumerate}[leftmargin=*]
    \item Add a \texttt{Measurement} type carrying label, \texttt{Point3},
          confidence, uncertainty, and timestamp.
    \item Add an overhead-camera calibration representation for image-to-table
          homography.
    \item Add a perception trait that can be backed by fiducials first and ML
          later.
    \item Add filtering and measurement gating before control.
    \item Add a task state machine for approach, grasp, lift, transfer, and
          release.
    \item Add real serial transport only after field measurements confirm board
          protocol, baud rate, IDs, and safe limits.
    \item Replace the crude gain mapping with measured Jacobian control or full
          inverse kinematics once link geometry is known.
\end{enumerate}

\section{Review Questions}

\begin{itemize}[leftmargin=*]
    \item Is the first demo explicitly tabletop pick-and-place?
    \item Should initial feedback be one overhead camera only, or overhead plus
          side camera?
    \item Are fiducials acceptable during calibration and data collection?
    \item Which objects should the first ML detector localize?
    \item What maximum motion speed and positional tolerance feel safe for the
          physical kit?
\end{itemize}

\end{document}
